\part{G\'eom\'etrie}

%% ════════════════════════════════════════════════════════════════
\chapter{Vecteurs du plan}
\begin{prerequis}
Lire les coordonnées d'un point; reconnaître parallélogramme, translation,
symétrie centrale et milieu d'un segment.
\end{prerequis}
\begin{automatismes}
Dans un repère, avec $A(1;2)$ et $B(4;-1)$, calculer les coordonnées de
$\vv{AB}$ et du milieu de $[AB]$; placer le point $C$ tel que
$\vv{AC}=2\vv{AB}$.
\end{automatismes}
\begin{activite}[Activité d'entrée -- Décrire un déplacement]
Sur un plan quadrillé, plusieurs chemins conduisent du point $A$ au point $B$.
Coder chacun par ses déplacements horizontal et vertical. Identifier ce qui
reste invariant malgré le chemin choisi, puis proposer une règle d'addition de
deux déplacements successifs.
\end{activite}


%% ════════════════════════════════════════════════════════════════

\begin{anecdote}[Hermann Grassmann et l'alg\`ebre vectorielle]
En $1844$, Hermann Grassmann publia \textit{Die Ausdehnungslehre}, introduisant
les espaces vectoriels en toute g\'en\'eralit\'e --- une id\'ee tellement en avance
sur son temps que peu de contemporains la comprirent. Aujourd'hui, les vecteurs
sont au c\oe ur de la physique, de l'infographie et de l'intelligence artificielle.
\end{anecdote}

\begin{objectifs}
\begin{itemize}
  \item Conna\^itre les caract\'eristiques d'un vecteur (direction, sens, norme).
  \item Ma\^itriser l'addition vectorielle et la multiplication par un scalaire.
  \item Calculer et utiliser le produit scalaire.
  \item D\'emontrer colin\'earit\'e et orthogonalit\'e.
\end{itemize}
\end{objectifs}

\section{G\'en\'eralit\'es}

\begin{defn}[Vecteur]
Un \textbf{vecteur} $\vv{AB}$ est d\'efini par sa \textbf{direction},
son \textbf{sens} et sa \textbf{norme} $\norme{\vv{AB}} = AB$.
Deux vecteurs sont \textbf{\'egaux} s'ils ont m\^{e}me direction, sens et norme.
\end{defn}

\begin{prop}[Milieu]
$I$ est le milieu de $[AB]$ si et seulement si
$\vv{IA} + \vv{IB} = \vv{0}$.
\end{prop}

\section{Produit scalaire}

\begin{defn}[Produit scalaire]
Soient $\vec{u}=(x_1,y_1)$ et $\vec{v}=(x_2,y_2)$ :
$$\vec{u} \cdot \vec{v} = x_1 x_2 + y_1 y_2
  = \norme{\vec{u}}\,\norme{\vec{v}}\cos\theta,$$
o\`u $\theta$ est l'angle entre $\vec{u}$ et $\vec{v}$.
\end{defn}

\begin{prop}[Orthogonalit\'e et colin\'earit\'e]
\begin{itemize}
  \item $\vec{u} \perp \vec{v} \Leftrightarrow \vec{u}\cdot\vec{v}=0$.
  \item $\vec{u} \parallel \vec{v} \Leftrightarrow x_1 y_2 - x_2 y_1 = 0$.
\end{itemize}
\end{prop}

\begin{figure}[H]
\centering
\begin{tikzpicture}[scale=1.15]
  \draw[-Stealth,BN,very thick] (0,0)--(3,1)
    node[midway,above,font=\small,text=BN] {$\vec{u}$};
  \draw[-Stealth,OR,very thick] (0,0)--(1.2,2.4)
    node[midway,left,font=\small,text=OR] {$\vec{v}$};
  \draw[-Stealth,VE,very thick] (0,0)--(4.2,3.4)
    node[midway,right,font=\small,text=VE] {$\vec{u}+\vec{v}$};
  \draw[dashed,BN!40] (3,1)--(4.2,3.4);
  \draw[dashed,OR!40] (1.2,2.4)--(4.2,3.4);
  \fill[BN] (0,0) circle (2pt) node[below left,font=\small] {$O$};
  \draw[GM,-Stealth] (0.65,0.22) arc[start angle=18,end angle=63,radius=0.7cm];
  \node[GM,font=\small] at (0.85,0.65) {$\theta$};
\end{tikzpicture}
\caption{Addition de vecteurs et angle form\'e.}
\end{figure}


\section{Op\'erations sur les vecteurs}

\begin{defn}[Addition et multiplication par un scalaire]
Soient $\vec{u}=(x_1,y_1)$, $\vec{v}=(x_2,y_2)$ et $\lambda\in\R$.
$$\vec{u}+\vec{v} = (x_1+x_2,\, y_1+y_2),\qquad
  \lambda\vec{u} = (\lambda x_1,\, \lambda y_1).$$
Le vecteur $-\vec{u} = (-x_1,-y_1)$ est l'\textbf{oppos\'e} de $\vec{u}$.
\end{defn}

\begin{prop}[Propri\'et\'es de l'addition vectorielle]
\begin{itemize}
  \item \textbf{Commutativit\'e~:} $\vec{u}+\vec{v}=\vec{v}+\vec{u}$.
  \item \textbf{Associativit\'e~:} $(\vec{u}+\vec{v})+\vec{w}=\vec{u}+(\vec{v}+\vec{w})$.
  \item \textbf{\'El\'ement neutre~:} $\vec{u}+\vec{0}=\vec{u}$.
  \item \textbf{Relation de Chasles~:} $\vv{AB}+\vv{BC}=\vv{AC}$.
\end{itemize}
\end{prop}

\begin{methode}[Coordonn\'ees d'un vecteur]
Si $A(x_A,y_A)$ et $B(x_B,y_B)$, alors~:
$$\vv{AB} = (x_B-x_A,\; y_B-y_A),\qquad
  \norme{\vv{AB}} = \sqrt{(x_B-x_A)^2+(y_B-y_A)^2}.$$
\end{methode}

\begin{defn}[Vecteurs colin\'eaires]
\index{colinéarité}
$\vec{u}=(x_1,y_1)$ et $\vec{v}=(x_2,y_2)$ sont \textbf{colin\'eaires}
(m\^eme direction) si et seulement si~:
$$x_1 y_2 - x_2 y_1 = 0.$$
Cette expression s'appelle le \textbf{d\'eterminant} du couple $(\vec{u},\vec{v})$.
\end{defn}

\begin{figure}[H]
\centering
\begin{tikzpicture}[scale=1.2]
  % Relation de Chasles: AB + BC = AC
  \coordinate (A) at (0,0);
  \coordinate (B) at (3,1.2);
  \coordinate (C) at (4.5,-0.5);
  % vectors
  \draw[-Stealth,BN,very thick] (A)--(B)
    node[midway,above left,font=\small] {$\vv{AB}$};
  \draw[-Stealth,OR,very thick] (B)--(C)
    node[midway,above right,font=\small] {$\vv{BC}$};
  \draw[-Stealth,VE,very thick,shorten >=1pt] (A)--(C)
    node[midway,below,font=\small] {$\vv{AC}=\vv{AB}+\vv{BC}$};
  % points
  \foreach \pt/\lbl/\pos in {A/{$A$}/below left, B/{$B$}/above, C/{$C$}/below right}{
    \fill[BN] (\pt) circle (2pt);
    \node[\pos,font=\small] at (\pt) {\lbl};
  }
  % second example: addition with parallelogram rule
  \begin{scope}[xshift=7cm]
    \coordinate (O) at (0,0);
    \coordinate (U) at (2.5,0.8);
    \coordinate (V) at (0.8,2);
    \coordinate (W) at (3.3,2.8);
    \draw[-Stealth,BN,very thick] (O)--(U)
      node[midway,below right,font=\small] {$\vec{u}$};
    \draw[-Stealth,OR,very thick] (O)--(V)
      node[midway,left,font=\small] {$\vec{v}$};
    \draw[-Stealth,VE,very thick] (O)--(W)
      node[midway,above left,font=\small] {$\vec{u}+\vec{v}$};
    \draw[dashed,BN!40] (U)--(W);
    \draw[dashed,OR!40] (V)--(W);
    \fill[BN] (O) circle (2pt);
    \node[below left,font=\small] at (O) {$O$};
  \end{scope}
  % labels below
  \node[font=\small,text=BN!70] at (2.2,-1.0)
    {Relation de Chasles};
  \node[font=\small,text=BN!70] at (9.4,-1.0)
    {R\`egle du parall\'elogramme};
\end{tikzpicture}
\caption{Addition de vecteurs~: relation de Chasles (gauche)
et r\`egle du parall\'elogramme (droite).}
\end{figure}

\begin{prop}[Isobarycentre et milieu]
Le milieu $I$ de $[AB]$ a pour coordonn\'ees~:
$I\!\left(\dfrac{x_A+x_B}{2},\; \dfrac{y_A+y_B}{2}\right).$

\noindent
Le centre de gravit\'e $G$ du triangle $ABC$ v\'erifie~:
$\vv{GA}+\vv{GB}+\vv{GC}=\vec{0}$, et ses coordonn\'ees sont
$G\!\left(\dfrac{x_A+x_B+x_C}{3},\; \dfrac{y_A+y_B+y_C}{3}\right).$
\end{prop}

\begin{figure}[H]
\centering
\begin{tikzpicture}[scale=1.0]
  %% ── Gauche : règle du parallélogramme ─────────────────────
  \draw[-Stealth,BN,very thick] (0,0)--(3,1)
    node[midway,above,font=\small,text=BN] {$\vec{u}$};
  \draw[-Stealth,OR,very thick] (0,0)--(1,2)
    node[midway,left,font=\small,text=OR] {$\vec{v}$};
  \draw[-Stealth,VE,very thick] (0,0)--(4,3)
    node[pos=0.65,above right,font=\small,text=VE] {$\vec{u}+\vec{v}$};
  \draw[dashed,BN!40] (3,1)--(4,3);
  \draw[dashed,OR!40] (1,2)--(4,3);
  \fill[BN] (0,0) circle (1.5pt);
  \node[below left,font=\small] at (0,0) {$O$};
  \node[font=\small,text=BN!60] at (2,-0.5) {Parallélogramme};
  %% ── Droite : multiplication par un scalaire ──────────────
  % séparateur vertical
  \draw[BN!20,dashed] (5.0,-0.8)--(5.0,3.5);
  % vecteur u entier
  \draw[-Stealth,BN,thick] (5.8,0)--(5.8+2.8,0.93)
    node[midway,below,font=\small,text=BN] {$\vec{u}$};
  % demi-vecteur lambda*u, décalé au-dessus
  \draw[-Stealth,TQ,thick] (5.8,1.4)--(5.8+1.4,1.4+0.47)
    node[midway,above,font=\small,text=TQ] {$\tfrac{1}{2}\vec{u}$};
  % vecteur -u, en dessous
  \draw[-Stealth,BO,thick] (5.8,-0.8)--(5.8-2.8,-0.8-0.93)
    node[midway,below,font=\small,text=BO] {$-\vec{u}$};
  \fill[BN] (0,0) circle (1.5pt);
  \node[font=\small,text=BN!60] at (7.8,-1.5) {Mult. par un scalaire};
\end{tikzpicture}
\caption{Gauche~: addition de vecteurs (r\`egle du parall\'elogramme).
Droite~: multiplication de $\vec{u}$ par un scalaire~: $\frac{1}{2}\vec{u}$
(demi-vecteur, bleu)~; $-\vec{u}$ (oppos\'e, rouge).}
\end{figure}

\section{Rep\`ere et coordonn\'ees}

\begin{defn}[Rep\`ere orthonorm\'e]
\index{repère orthonormé}
Un \textbf{rep\`ere orthonorm\'e} $(O, \vec{\imath}, \vec{\jmath})$ est form\'e d'un
point $O$ (origine) et de deux vecteurs perpendiculaires de norme $1$.
Tout vecteur $\vec{u}$ s'\'ecrit $\vec{u} = x\vec{\imath} + y\vec{\jmath}$~;
$(x,y)$ sont les \textbf{coordonn\'ees} de $\vec{u}$.
\end{defn}

\begin{methode}[D\'emontrer que des points sont align\'es]
\index{alignement}
$A$, $B$, $C$ sont align\'es ssi $\vv{AB}$ et $\vv{AC}$ sont colin\'eaires, c'est-\`a-dire~:
$$\det(\vv{AB}, \vv{AC})
  = (x_B-x_A)(y_C-y_A) - (x_C-x_A)(y_B-y_A) = 0.$$
\end{methode}

\begin{methode}[D\'emontrer qu'un quadrilat\`ere est un parall\'elogramme]
$ABCD$ est un parall\'elogramme ssi $\vv{AB} = \vv{DC}$
(ou, de mani\`ere \'equivalente, si les diagonales se coupent en leur milieu).
\end{methode}

\section{Exercices}

\subsection*{\diff{1}\quad Niveau 1 \textemdash\ Calcul de base}

\begin{exo}[Coordonn\'ees et norme \corrige]{1}
On donne $A(1,2)$, $B(4,6)$, $C(-1,0)$, $D(3,-2)$.
\begin{enumerate}
  \item Calculer $\vv{AB}$, $\vv{BA}$, $\vv{AC}$, $\vv{CD}$.
  \item Calculer $\norme{\vv{AB}}$, $\norme{\vv{AC}}$, $\norme{\vv{CD}}$.
  \item Calculer les coordonn\'ees du milieu $M$ de $[AB]$.
  \item Calculer le centre de gravit\'e $G$ du triangle $ABC$.
\end{enumerate}
\end{exo}

\begin{exo}[Op\'erations vectorielles \corrige]{1}
On pose $\vec{u}=(2,-1)$ et $\vec{v}=(-1,3)$.
\begin{enumerate}
  \item Calculer $\vec{u}+\vec{v}$, $\vec{u}-\vec{v}$,
        $2\vec{u}+3\vec{v}$, $\vec{u}-2\vec{v}$.
  \item Calculer $\norme{\vec{u}}$, $\norme{\vec{v}}$,
        $\norme{\vec{u}+\vec{v}}$.
  \item V\'erifier l'in\'egalit\'e triangulaire~:
        $\norme{\vec{u}+\vec{v}} \leq \norme{\vec{u}}+\norme{\vec{v}}$.
\end{enumerate}
\end{exo}

\begin{exo}[Produits scalaires \corrige]{1}
Calculer le produit scalaire, puis l'angle entre les vecteurs~:
\begin{multicols}{2}
\begin{enumerate}
  \item $\vec{u}=(3,4)$, $\vec{v}=(1,-2)$
  \item $\vec{u}=(2,1)$, $\vec{v}=(-1,2)$
  \item $\vec{u}=(5,0)$, $\vec{v}=(0,7)$
  \item $\vec{u}=(1,1)$, $\vec{v}=(1,-1)$
  \item $\vec{u}=(3,0)$, $\vec{v}=(0,-4)$
  \item $\vec{u}=(1,\sqrt{3})$, $\vec{v}=(\sqrt{3},1)$
\end{enumerate}
\end{multicols}
\end{exo}

\begin{exo}[Parall\'elogramme \corrige]{1}
V\'erifier que $ABCD$ est un parall\'elogramme avec~:
\begin{enumerate}
  \item $A(1,1)$, $B(4,3)$, $C(3,0)$, $D(0,-2)$.
  \item $A(-1,2)$, $B(3,4)$, $C(5,1)$, $D(1,-1)$.
\end{enumerate}
Pour chacun, calculer les longueurs des c\^ot\'es et des diagonales.
\end{exo}

\begin{exo}[Orthogonalit\'e \corrige]{1}
D\'eterminer la valeur de $k$ telle que $\vec{u}$ et $\vec{v}$ soient orthogonaux~:
\begin{multicols}{2}
\begin{enumerate}
  \item $\vec{u}=(k,3)$, $\vec{v}=(2,k-1)$
  \item $\vec{u}=(k,2k)$, $\vec{v}=(4,k)$
  \item $\vec{u}=(1,k)$, $\vec{v}=(k,9)$
  \item $\vec{u}=(k+1,2)$, $\vec{v}=(3,k-1)$
\end{enumerate}
\end{multicols}
\end{exo}

\begin{exo}[Colin\'earit\'e et alignement \corrige]{1}
\begin{enumerate}
  \item D\'eterminer si les vecteurs sont colin\'eaires~:
        $\vec{u}=(2,4)$ et $\vec{v}=(1,2)$~;
        $\vec{u}=(3,1)$ et $\vec{v}=(6,3)$~;
        $\vec{u}=(1,-2)$ et $\vec{v}=(2,4)$.
  \item Les points $A(1,2)$, $B(3,6)$, $C(5,10)$ sont-ils align\'es~?
  \item Les points $P(0,1)$, $Q(2,4)$, $R(4,8)$ sont-ils align\'es~?
  \item D\'eterminer $k$ pour que $A(1,0)$, $B(k,2)$, $C(3,4)$ soient align\'es.
\end{enumerate}
\end{exo}

\begin{exo}[Relation de Chasles \corrige]{1}
\index{relation de Chasles}
En utilisant la relation de Chasles, simplifier~:
\begin{multicols}{2}
\begin{enumerate}
  \item $\vv{AB}+\vv{BC}$
  \item $\vv{AB}+\vv{BC}+\vv{CA}$
  \item $\vv{MA}+\vv{AB}+\vv{BM}$
  \item $\vv{AB}-\vv{AC}$
  \item $\vv{OA}+\vv{OB}+\vv{OC}$ quand $G=O$ est le centre de gravité de $ABC$.
  \item $2\vv{GA}+2\vv{GB}+2\vv{GC}$
\end{enumerate}
\end{multicols}
\end{exo}

\begin{exo}[Norme et distance \corrige]{1}
On donne $A(-2,3)$, $B(4,-1)$, $C(1,5)$.
\begin{enumerate}
  \item Calculer les longueurs $AB$, $BC$, $CA$.
  \item Le triangle $ABC$ est-il \'equilat\'eral~? Isoc\`ele~?
  \item Calculer le p\'erim\`etre du triangle.
  \item Trouver les coordonn\'ees du milieu de chaque c\^ot\'e.
\end{enumerate}
\end{exo}

\begin{exo}[Vecteur directeur d'une droite \corrige]{1}
\index{vecteur directeur}
Un vecteur directeur d'une droite est un vecteur non nul parall\`ele
\`a cette droite.
\begin{enumerate}
  \item La droite $y = 2x - 1$ a pour vecteur directeur $(1, 2)$.
        V\'erifier.
  \item Donner un vecteur directeur de la droite $y = -3x + 5$.
  \item Donner un vecteur directeur de la droite $2x - y + 1 = 0$.
  \item Deux droites de vecteurs directeurs $\vec{u}$ et $\vec{v}$
        sont parall\`eles ssi $\vec{u}$ et $\vec{v}$ sont colin\'eaires.
        D\'eterminer si les droites $y=2x+1$ et $4x-2y-3=0$ sont parall\`eles.
\end{enumerate}
\end{exo}

\begin{exo}[Vecteur normal \`a une droite \corrige]{1}
\index{vecteur normal}
Un vecteur \textbf{normal} \`a une droite lui est perpendiculaire.
\begin{enumerate}
  \item La droite $ax+by+c=0$ a pour vecteur normal $(a,b)$.
        V\'erifier pour $2x-3y+1=0$.
  \item Donner un vecteur normal \`a $y = x+1$.
  \item Deux droites sont perpendiculaires ssi leurs vecteurs normaux sont
        perpendiculaires. D\'eterminer si $x+2y=0$ et $2x-y+3=0$ sont perpendiculaires.
  \item \'Ecrire l'\'equation d'une droite passant par $A(1,2)$ et
        de vecteur normal $(3,-1)$.
\end{enumerate}
\end{exo}

\begin{exo}[Formulaires vectoriels \corrige]{1}
Calculer, sachant que $\norme{\vec{u}}=3$, $\norme{\vec{v}}=4$
et $\vec{u}\cdot\vec{v}=6$~:
\begin{multicols}{2}
\begin{enumerate}
  \item $\norme{\vec{u}+\vec{v}}^2$
  \item $\norme{\vec{u}-\vec{v}}^2$
  \item $(\vec{u}+\vec{v})\cdot(\vec{u}-\vec{v})$
  \item $\norme{2\vec{u}-\vec{v}}^2$
  \item L'angle entre $\vec{u}$ et $\vec{v}$
  \item $\norme{\vec{u}+\vec{v}}$
\end{enumerate}
\end{multicols}
\end{exo}

\begin{exo}[Construction de figures par vecteurs \corrige]{1}
On donne $A(0,0)$, $B(4,0)$, $C(2,3)$.
\begin{enumerate}
  \item Calculer $\vv{AB}$, $\vv{BC}$, $\vv{CA}$.
  \item Calculer $\vv{AB}\cdot\vv{BC}$.
        Le triangle est-il rectangle en $B$~?
  \item Trouver le point $D$ tel que $ABDC$ soit un parall\'elogramme.
  \item Trouver le point $M$ milieu de $[BD]$.
        V\'erifier qu'il co\"incide avec le milieu de $[AC]$.
\end{enumerate}
\end{exo}

\begin{exo}[Vecteur et valeur param\'etrique \corrige]{1}
La droite $(d)$ passe par $A(1,2)$ et a pour vecteur directeur $\vec{u}=(3,-1)$.
\begin{enumerate}
  \item \'Ecrire les \'equations param\'etriques de $(d)$~:
        $\begin{cases}x=1+3t\\y=2-t\end{cases}$, $t\in\R$.
  \item Pour $t=0$, $t=1$, $t=-1$~: calculer les points correspondants.
  \item Le point $P(7,0)$ appartient-il \`a $(d)$~?
  \item \'Ecrire l'\'equation cart\'esienne de $(d)$.
\end{enumerate}
\end{exo}

\begin{exo}[Point image par translation \corrige]{1}
Une translation de vecteur $\vec{t}=(2,-3)$ transforme $M$ en $M'$.
\begin{enumerate}
  \item Si $M(1,4)$, calculer les coordonn\'ees de $M'$.
  \item Si $M'(5,1)$, calculer les coordonn\'ees de $M$.
  \item Quel est l'image du segment $[AB]$ avec $A(0,0)$, $B(3,2)$~?
  \item Quel est l'image du triangle $OAB$ ($O$ origine)~?
\end{enumerate}
\end{exo}

\begin{exo}[Homoth\'etie et vecteur \corrige]{1}
\index{homothétie}
L'homoth\'etie de centre $O$ et de rapport $k$ transforme $\vv{OM}$
en $\vv{OM'} = k\vv{OM}$.
\begin{enumerate}
  \item Avec $O(0,0)$, $k=2$~: transformer $A(1,3)$, $B(-1,2)$, $C(2,-1)$.
  \item Avec $O(1,1)$, $k=3$~: transformer $A(2,3)$.
        \emph{(Indication~: $\vv{OA'} = 3\vv{OA}$~; calculer $A'$.)}
  \item Quelle est l'image du segment $[AB]$ ($A(1,0)$, $B(3,2)$) par
        l'homoth\'etie de centre $O(0,0)$ et rapport $-1$~?
\end{enumerate}
\end{exo}

\subsection*{\diff{2}\quad Niveau 2 \textemdash\ Applications et d\'emonstrations}

\begin{exo}[Angle et produit scalaire \corrige]{2}
On donne $A(2,3)$, $B(5,1)$, $C(4,6)$.
\begin{enumerate}
  \item Calculer $\vv{AB}\cdot\vv{AC}$.
  \item Le triangle $ABC$ est-il rectangle en $A$~? En $B$~? En $C$~?
  \item Calculer l'angle $\hat{A}$ du triangle \`a $1^{\circ}$ pr\`es.
  \item Calculer l'aire du triangle.
\end{enumerate}
\end{exo}

\begin{exo}[D\'emontrer des propri\'et\'es g\'eom\'etriques \corrige]{2}
D\'emontrer chacune des propri\'et\'es suivantes par le calcul vectoriel.
\begin{enumerate}
  \item Les diagonales d'un parall\'elogramme se coupent en leur milieu.
  \item Dans un triangle $ABC$, le milieu $M$ de $[BC]$ v\'erifie
        $\vv{AM} = \frac{1}{2}(\vv{AB}+\vv{AC})$.
  \item Si $\vv{OA}+\vv{OB}+\vv{OC} = \vec{0}$, alors $O$ est le centre de gravité de $ABC$.
\end{enumerate}
\end{exo}

\begin{exo}[Formule de polarisation \corrige]{2}
On admet que~:
$\vec{u}\cdot\vec{v} = \dfrac{1}{2}\bigl(\norme{\vec{u}}^2+\norme{\vec{v}}^2
-\norme{\vec{u}-\vec{v}}^2\bigr).$
\begin{enumerate}
  \item V\'erifier cette formule en coordonn\'ees pour $\vec{u}=(a,b)$
        et $\vec{v}=(c,d)$.
  \item Appliquer \`a $A(1,2)$, $B(4,6)$, $C(0,-1)$ pour calculer
        $\vv{AB}\cdot\vv{AC}$ sans coordonn\'ees, \`a partir des distances.
  \item En d\'eduire l'angle $\hat{A}$ du triangle.
\end{enumerate}
\end{exo}

\begin{exo}[Formule du parall\'elogramme \corrige]{2}
D\'emontrer la formule~:
$$\norme{\vec{u}+\vec{v}}^2 + \norme{\vec{u}-\vec{v}}^2
  = 2\norme{\vec{u}}^2 + 2\norme{\vec{v}}^2.$$
\begin{enumerate}
  \item D\'evelopper $\norme{\vec{u}+\vec{v}}^2$ et $\norme{\vec{u}-\vec{v}}^2$
        en utilisant $\norme{\vec{w}}^2=\vec{w}\cdot\vec{w}$.
  \item Additionner et conclure.
  \item Interpr\'etation g\'eom\'etrique~: dans le parall\'elogramme $OABC$
        (o\`u $\vec{u}=\vv{OA}$ et $\vec{v}=\vv{OC}$), quelle est la signification
        de chaque terme~?
\end{enumerate}
\end{exo}

\begin{exo}[Alignement et colin\'earit\'e \corrige]{2}
\begin{enumerate}
  \item D\'emontrer que $A(1,2)$, $B(3,6)$, $C(5,10)$ sont align\'es.
  \item Trouver le point $D$ sur la droite $(AB)$ tel que $AD = 2\cdot AB$.
  \item D\'eterminer si le quadrilat\`ere $ABCD$ avec $A(0,0)$, $B(3,1)$,
        $C(5,4)$, $D(2,3)$ est un parall\'elogramme, un losange, ou ni l'un ni l'autre.
\end{enumerate}
\end{exo}

\begin{exo}[Vecteurs et g\'eom\'etrie classique \corrige]{2}
On consid\`ere un triangle $ABC$ avec $A(0,0)$, $B(6,0)$, $C(2,4)$.
\begin{enumerate}
  \item Calculer les coordonn\'ees des milieux $I_{AB}$, $I_{BC}$, $I_{CA}$
        de chaque c\^ot\'e.
  \item Montrer que les vecteurs $\vv{AI_{BC}}$ et $\vv{BI_{CA}}$
        se coupent en un point $G$~: calculer ses coordonn\'ees.
  \item V\'erifier que $G$ est bien le centre de gravité de $ABC$.
  \item Calculer $AG$ et $GI_{BC}$. Quel est le rapport $AG/GI_{BC}$~?
\end{enumerate}
\end{exo}

\begin{exo}[Projection orth. et distance \corrige]{2}
\index{projection orthogonale}
La projection orthogonale de $A$ sur une droite $(d)$ est le pied $H$
de la perpendiculaire men\'ee de $A$ \`a $(d)$.
\begin{enumerate}
  \item La droite $(d)$ a pour \'equation $2x+y-5=0$ et $A=(3,1)$.
        Trouver les coordonn\'ees de $H$ en \'ecrivant que $H\in(d)$
        et $\vv{AH}\cdot\vec{n}=0$ o\`u $\vec{n}$ est le vecteur normal \`a $(d)$.
  \item Calculer la distance $AH$.
  \item G\'en\'eraliser~: montrer que $d(A,(d)) = \dfrac{|ax_A+by_A+c|}{\sqrt{a^2+b^2}}$.
\end{enumerate}
\end{exo}

\begin{exo}[Vecteurs et calcul d'angle \corrige]{2}
Dans un triangle $ABC$ avec $A(2,0)$, $B(6,2)$, $C(4,6)$~:
\begin{enumerate}
  \item Calculer les trois angles du triangle.
  \item V\'erifier que leur somme vaut bien $180^{\circ}$.
  \item Calculer les aires par la m\'ethode $\frac{1}{2}\norme{\vec{u}}\norme{\vec{v}}|\sin\theta|$.
\end{enumerate}
\end{exo}

\begin{exo}[Coordonn\'ees et \'equations de droites \corrige]{2}
\begin{enumerate}
  \item \'Ecrire l'\'equation de la droite passant par $A(2,-1)$ et
        perpendiculaire \`a $\vec{n}=(3,4)$.
  \item \'Ecrire l'\'equation de la droite passant par $A(1,2)$ et
        $B(3,-1)$ sous forme $ax+by+c=0$.
  \item D\'eterminer la droite m\'ediatrice du segment $[AB]$
        avec $A(2,6)$, $B(8,2)$.
        \emph{(Elle passe par le milieu de $[AB]$ et lui est perpendiculaire.)}
  \item V\'erifier que le centre du cercle circumscrit \`a un triangle
        est l'intersection des trois m\'ediatrices.
\end{enumerate}
\end{exo}

\begin{exo}[Milieu et centre de gravité : applications \corrige]{2}
$A(2,0)$, $B(8,4)$, $C(4,10)$.
\begin{enumerate}
  \item Calculer le milieu de chaque c\^ot\'e.
  \item Calculer le centre de gravité $G$.
  \item V\'erifier que $\vv{GA}+\vv{GB}+\vv{GC}=\vec{0}$.
  \item Montrer que les m\'edianes se coupent en $G$.
\end{enumerate}
\end{exo}

\begin{exo}[\'Equation vectorielle \corrige]{2}
Trouver le point $M(x,y)$ v\'erifiant~:
\begin{enumerate}
  \item $\vv{OM} = 2\vv{OA} - \vv{OB}$ avec $A(3,1)$, $B(1,5)$.
  \item $\vv{AM} = 3\vv{AB}$ avec $A(1,2)$, $B(4,5)$.
  \item $\vv{MA} + \vv{MB} = \vec{0}$~: que repr\'esente $M$~?
  \item $\norme{\vv{OM}} = 5$ et $\vv{OM} = t(1,2)$ pour un r\'eel $t>0$.
\end{enumerate}
\end{exo}

\begin{exo}[Sym\'etrie centrale par vecteurs \corrige]{2}
Le sym\'etrique $M'$ de $M$ par rapport \`a un centre $C$ v\'erifie
$\vv{CM'} = -\vv{CM}$.
\begin{enumerate}
  \item Trouver le sym\'etrique de $A(3,5)$ par rapport \`a $C(1,2)$.
  \item Trouver le sym\'etrique de $B(-2,4)$ par rapport \`a $O(0,0)$.
  \item Montrer que $\vv{CM'}=-\vv{CM}$ implique que $C$ est le milieu de $[MM']$.
  \item La sym\'etrie centrale conserve-t-elle les distances~? V\'erifier.
\end{enumerate}
\end{exo}

\begin{exo}[Produit scalaire et angle dans un triangle \corrige]{2}
Pour $A(1,0)$, $B(4,0)$, $C(2,3)$~:
\begin{enumerate}
  \item Calculer tous les angles du triangle avec le produit scalaire.
  \item D\'eterminer s'il est acutangle, rectangle ou obtusangle.
  \item Calculer l'aire du triangle par deux m\'ethodes.
\end{enumerate}
\end{exo}

\begin{exo}[Condition d'alignement \corrige]{1}
D\'eterminer si les points sont align\'es~:
\begin{enumerate}
  \item $A(0,0)$, $B(2,3)$, $C(4,6)$
  \item $A(1,2)$, $B(3,5)$, $C(7,11)$
  \item $A(-1,4)$, $B(2,1)$, $C(5,-3)$
  \item $A(0,3)$, $B(4,1)$, $C(8,0)$
\end{enumerate}
Pour (d), trouver $k$ tel que $A(0,3)$, $B(4,k)$, $C(8,0)$ soient align\'es.
\end{exo}

\begin{exo}[Quadrilat\`ere et vecteurs \corrige]{1}
Pour $A(0,0)$, $B(4,1)$, $C(5,4)$, $D(1,3)$~:
\begin{enumerate}
  \item Calculer $\vv{AB}$, $\vv{BC}$, $\vv{CD}$, $\vv{DA}$.
  \item V\'erifier que $\vv{AB}=\vv{DC}$. Quelle est la nature de $ABCD$~?
  \item Calculer les longueurs des c\^ot\'es et des diagonales.
  \item $ABCD$ est-il un rectangle~? Un losange~?
\end{enumerate}
\end{exo}

\begin{exo}[\logicoalgo\ Algorithme~: tester l'alignement de trois points \corrige]{2}
Trois points $A(x_A,y_A)$, $B(x_B,y_B)$, $C(x_C,y_C)$ sont align\'es
si et seulement si $\det(\vv{AB},\vv{AC}) = (x_B-x_A)(y_C-y_A)-(x_C-x_A)(y_B-y_A) = 0$.

\begin{algorithm}[H]
\caption{Test d'alignement de trois points}
\KwIn{les coordonn\'ees de $A$, $B$, $C$}
\KwOut{\Vrai\ si align\'es, \Faux\ sinon}
$d \leftarrow (x_B-x_A)\times(y_C-y_A) - (x_C-x_A)\times(y_B-y_A)$\;
\Si{$d = 0$}{\Afficher \og Align\'es\fg\;}
\Sinon{\Afficher \og Non align\'es\fg\;}
\end{algorithm}

\begin{enumerate}
  \item La condition $d=0$ est-elle n\'ecessaire \emph{et} suffisante~?
        \'Enoncer cela sous forme d'\'equivalence logique.
  \item Tester avec $A(0,0)$, $B(2,3)$, $C(4,6)$~;
        puis avec $A(1,1)$, $B(2,3)$, $C(4,6)$.
  \item La condition ``$|d| < \varepsilon$'' (avec $\varepsilon=10^{-9}$)
        est-elle plus prudente que $d=0$~?
        Pourquoi dans un programme~?
  \item \'Ecrire une version Python \computer\ prenant les coordonn\'ees
        en entr\'ee et retournant \texttt{True} ou \texttt{False}.
\end{enumerate}
\end{exo}

\subsection*{\diff{3}\quad Niveau 3 \textemdash\ D\'emonstrations vectorielles}

\begin{exo}[Losange et produit scalaire \corrige]{3}
D\'emontrer, \`a l'aide des vecteurs, que les diagonales d'un losange
sont perpendiculaires.
Soit $ABCD$ un losange (quatre c\^ot\'es \'egaux~: $AB=BC=CD=DA=a$).
\begin{enumerate}
  \item Exprimer $\vv{AC}$ et $\vv{BD}$ en fonction de $\vv{AB}$ et $\vv{AD}$.
  \item Calculer $\vv{AC}\cdot\vv{BD}$.
  \item Conclure.
\end{enumerate}
\end{exo}

\begin{exo}[Th\'eor\`eme de la m\'ediane \corrige]{3}
\index{théorème de la médiane}
Dans tout triangle $ABC$, si $M$ est le milieu de $[BC]$~:
$$AB^2 + AC^2 = 2AM^2 + 2BM^2.$$
\begin{enumerate}
  \item Exprimer $\vv{AB}$ et $\vv{AC}$ en fonction de $\vv{AM}$
        et $\vv{BM}$ (avec $\vv{MB}=-\vv{BM}$).
  \item Calculer $AB^2+AC^2 = \norme{\vv{AB}}^2+\norme{\vv{AC}}^2$.
  \item Conclure.
  \item Application~: $A(1,1)$, $B(5,3)$, $C(-1,7)$.
        V\'erifier num\'eriquement.
\end{enumerate}
\end{exo}

\begin{exo}[Th\'eor\`eme de l'angle inscrit dans un demi-cercle \corrige]{3}
Soit un cercle de centre $O$ et un diam\`etre $[AB]$. Pour tout $M$
sur le cercle ($M\ne A$, $M\ne B$), montrer que $\widehat{AMB} = 90^{\circ}$.
\begin{enumerate}
  \item Poser $\vv{OA}=\vec{a}$ et $\vv{OM}=\vec{m}$
        avec $\norme{\vec{a}}=\norme{\vec{m}}=R$ (rayon).
  \item Exprimer $\vv{MA}$ et $\vv{MB}$ en fonction de $\vec{a}$ et $\vec{m}$.
  \item Calculer $\vv{MA}\cdot\vv{MB}$.
  \item Conclure.
\end{enumerate}
\end{exo}

\begin{exo}[Centre de gravité et vecteurs \corrige]{3}
\begin{enumerate}
  \item D\'emontrer que si $G$ est le centre de gravité de $ABC$, alors pour tout
        point $P$~: $\vv{PA}+\vv{PB}+\vv{PC} = 3\vv{PG}$.
  \item En d\'eduire~: $\vv{GA}+\vv{GB}+\vv{GC}=\vec{0}$.
  \item D\'emontrer que le centre de gravité partage chaque m\'ediane dans le rapport $2:1$.
        \emph{(Montrer que $G$ est sur la m\'ediane issue de $A$, et que $AG/GM_A=2$.)}
\end{enumerate}
\end{exo}

\begin{exo}[Loi des cosinus \corrige]{3}
\begin{enumerate}
  \item D\'emontrer, \`a l'aide du produit scalaire, la loi des cosinus~:
        $$a^2 = b^2 + c^2 - 2bc\cos\widehat{A},$$
        o\`u $a=BC$, $b=CA$, $c=AB$ et $\widehat{A}=\widehat{BAC}$.
        \emph{(Exprimer $\vv{BC}=\vv{AC}-\vv{AB}$ et calculer $\norme{\vv{BC}}^2$.)}
  \item V\'erifier avec $A(0,0)$, $B(4,0)$, $C(1,3)$.
  \item En quoi la loi des cosinus g\'en\'eralise-t-elle le th\'eor\`eme de Pythagore~?
\end{enumerate}
\end{exo}

\begin{exo}[Orthocentre d'un triangle \corrige]{3}
\index{orthocentre}
L'\textbf{orthocentre} $H$ d'un triangle $ABC$ est l'intersection
des trois hauteurs.
\begin{enumerate}
  \item Montrer que $H$ est caract\'eris\'e par~:
        $\vv{HA}\cdot\vv{BC}=0$ et $\vv{HB}\cdot\vv{AC}=0$.
  \item Pour $A(0,4)$, $B(-2,0)$, $C(3,0)$~:
        trouver l'\'equation de la hauteur issue de $A$,
        puis celle issue de $B$.
  \item Calculer les coordonn\'ees de $H$.
  \item V\'erifier que $\vv{HC}\cdot\vv{AB}=0$.
\end{enumerate}
\end{exo}

%─────────────────────────────────────────────────────────────────
\subsection*{Probl\`emes}
%─────────────────────────────────────────────────────────────────

\begin{probleme}[Navigation et cap \corrigeP]{1}
Un bateau quitte le port $O(0,0)$ et navigue successivement~:
$\vv{u_1}=(3,4)$ km, puis $\vv{u_2}=(-1,2)$ km, puis $\vv{u_3}=(2,-3)$ km.
\begin{enumerate}
  \item Calculer la position finale $P$ du bateau.
  \item Calculer la distance $OP$.
  \item Quel vecteur unique $\vec{t}$ m\`enerait directement de $O$ \`a $P$~?
  \item L'angle de cap (par rapport au Nord, i.e.\ l'axe $Oy$) vaut
        $\arctan(x_P/y_P)$ si $y_P>0$. Calculer cet angle.
\end{enumerate}
\end{probleme}

\begin{probleme}[G\'eom\'etrie du parall\'elogramme \corrigeP]{2}
Soit $ABCD$ un parall\'elogramme ($A(0,0)$, $B(4,0)$, $D(1,3)$).
\begin{enumerate}
  \item Calculer les coordonn\'ees de $C$.
  \item Calculer les longueurs des c\^ot\'es et des diagonales.
  \item D\'eterminer si $ABCD$ est un rectangle, un losange, ou les deux.
        \emph{(Rectangle~: diagonales \'egales~; losange~: c\^ot\'es \'egaux.)}
  \item Calculer l'aire de $ABCD$.
        \emph{(Aire~: $|x_B y_D - x_D y_B|$ pour un parall\'elogramme $OAD$.)}
\end{enumerate}
\end{probleme}

\begin{probleme}[Barycentre et partage d'un segment \corrigeP]{2}
Le point $G$ partage le segment $[AB]$ dans le rapport $k = AG/GB > 0$.
\begin{enumerate}
  \item Exprimer $\vv{OG}$ en fonction de $\vv{OA}$, $\vv{OB}$ et $k$.
  \item Pour $A(1,2)$, $B(5,6)$, trouver $G$ tel que $AG/GB = 3$.
  \item Trouver le point $G$ tel que $AG = 3\,GB$ avec $A(0,0)$, $B(8,4)$.
  \item G\'en\'eraliser~: si $G$ divise $[AB]$ dans le rapport $\lambda/(1-\lambda)$,
        montrer que $\vv{OG} = (1-\lambda)\vv{OA}+\lambda\vv{OB}$.
\end{enumerate}
\end{probleme}

\begin{probleme}[Vecteurs et physique~: forces \corrigeP]{1}
Deux forces $\vec{F_1}=(3,4)$ N et $\vec{F_2}=(-1,2)$ N s'appliquent
sur un objet.
\begin{enumerate}
  \item Calculer la force r\'esultante $\vec{R}=\vec{F_1}+\vec{F_2}$.
  \item Calculer $\norme{\vec{R}}$ (l'intensit\'e de la force r\'esultante).
  \item Une troisi\`eme force $\vec{F_3}$ s'ajoute pour que l'objet soit
        en \'equilibre ($\vec{R}+\vec{F_3}=\vec{0}$). Trouver $\vec{F_3}$.
  \item Calculer l'angle entre $\vec{F_1}$ et $\vec{F_2}$.
\end{enumerate}
\end{probleme}

\begin{probleme}[Cercle d\'efini par son diam\`etre \corrigeP]{2}
Soit $A(1,-1)$ et $B(5,3)$.
\begin{enumerate}
  \item Calculer le milieu $\Omega$ de $[AB]$ et le rayon $R = \dfrac{AB}{2}$.
  \item \'Ecrire l'\'equation du cercle de diam\`etre $[AB]$.
  \item Pour un point $M(x,y)$ quelconque~:
        montrer que $M$ est sur ce cercle ssi $\vv{MA}\cdot\vv{MB}=0$.
  \item V\'erifier avec $M(1,3)$ et $M(5,-1)$.
\end{enumerate}
\end{probleme}

\begin{probleme}[M\'ediane, hauteur, m\'ediatrice \corrigeP]{2}
Soit le triangle $A(0,6)$, $B(-4,0)$, $C(4,0)$.
\begin{enumerate}
  \item Calculer les coordonn\'ees du centre de gravité $G$,
        du circocentre $O'$ (intersection des m\'ediatrices)
        et de l'orthocentre $H$ (intersection des hauteurs).
  \item V\'erifier la \emph{droite d'Euler}~: $O'$, $G$, $H$ sont align\'es
        et $O'G = G H/2$.
  \item Le triangle est-il rectangle~? Isoc\`ele~? \'Equilat\'eral~?
\end{enumerate}
\end{probleme}

\begin{probleme}[Vecteurs et d\'eplacement \corrigeP]{1}
Sur une grille, un robot part de $A(0,0)$ et ex\'ecute les instructions~:
\begin{itemize}
  \item ``Nord'' = $+2$ en $y$~;
        ``Est'' = $+2$ en $x$~;
        ``Sud-Est'' = $(1,-1)$.
\end{itemize}
S\'equence~: Nord, Est, Nord, Sud-Est, Est.
\begin{enumerate}
  \item Calculer la position finale $B$.
  \item Quel est le vecteur $\vv{AB}$~?
  \item Quelle s\'equence d'exactement $3$ mouvements m\`enerait de $A$ \`a $B$~?
  \item \'Ecrire un algorithme Python \computer\ simulant cette trajectoire
        et affichant chaque position interm\'ediaire.
\end{enumerate}
\end{probleme}

\begin{probleme}[Vecteurs et r\'egression lin\'eaire \corrigeP]{2}
On dispose des points de mesures~:
$(1,2)$, $(2,3)$, $(3,5)$, $(4,4)$, $(5,6)$.
\begin{enumerate}
  \item Calculer la moyenne $\bar{x}$ et $\bar{y}$ des coordonn\'ees.
  \item Le vecteur $\vec{m}=(\bar{x},\bar{y})$ pointe vers le \og barycentre\fg\
        des donn\'ees. Calculer $\vec{m}$.
  \item Calculer $\sum(x_i-\bar{x})^2$ et $\sum(x_i-\bar{x})(y_i-\bar{y})$.
  \item La droite de r\'egression a pour coefficient directeur
        $a = \dfrac{\sum(x_i-\bar{x})(y_i-\bar{y})}{\sum(x_i-\bar{x})^2}$.
        Calculer $a$ et l'ordonn\'ee \`a l'origine $b=\bar{y}-a\bar{x}$.
\end{enumerate}
\end{probleme}

\begin{probleme}[Polygone r\'egulier \corrigeP]{3}
On construit un hexagone r\'egulier $ABCDEF$ de centre $O$ et de sommet
$A(2,0)$.
\begin{enumerate}
  \item Les six sommets sont obtenus en faisant tourner $\vv{OA}$
        de $60^{\circ}$ \`a chaque \'etape. Calculer les coordonn\'ees de
        $B$, $C$, $D$, $E$, $F$.
        \emph{(Rotation de $60^{\circ}$~: $(x,y)\mapsto(x\cos60^{\circ}-y\sin60^{\circ},\,x\sin60^{\circ}+y\cos60^{\circ})$.)}
  \item V\'erifier que $\vv{OA}+\vv{OB}+\vv{OC}+\vv{OD}+\vv{OE}+\vv{OF}=\vec{0}$.
  \item Calculer le produit scalaire $\vv{AB}\cdot\vv{BC}$.
  \item Calculer l'aire de l'hexagone.
\end{enumerate}
\end{probleme}

\begin{probleme}[Construction d'un point par vecteur \corrigeP]{1}
On donne $A(2,1)$, $B(5,4)$, $C(1,6)$.
\begin{enumerate}
  \item Trouver $D$ tel que $ABDC$ soit un parall\'elogramme.
  \item Trouver $E$ tel que $\vv{AE} = 2\vv{AB} + \vv{AC}$.
  \item Trouver $F$ milieu de $[BE]$.
  \item Définir le point $G$ par
        $\vv{OG}=\dfrac{\vv{OA}+\vv{OB}+\vv{OC}}{3}$, puis vérifier que $G$
        est le centre de gravité et
        que $\vv{GA}+\vv{GB}+\vv{GC}=\vec{0}$.
\end{enumerate}
\end{probleme}

\begin{probleme}[Vecteurs et r\`egle de Thal\`es \corrigeP]{2}
Soit $A(0,0)$, $B(6,0)$, $C(0,4)$. $M$ est le milieu de $[AB]$
et $N$ est le milieu de $[AC]$.
\begin{enumerate}
  \item Calculer les coordonn\'ees de $M$ et $N$.
  \item Calculer $\vv{MN}$ et $\vv{BC}$.
        Que remarque-t-on~?
  \item Calculer $MN$ et $BC$.
        Quel est le rapport $MN/BC$~?
  \item Montrer que $MN\parallel BC$ en utilisant la colin\'earit\'e des vecteurs.
\end{enumerate}
\end{probleme}

\begin{probleme}[Vecteurs et invariance par translation \corrigeP]{3}
Une translation de vecteur $\vec{t}=(a,b)$ transforme $M(x,y)$ en $M'(x+a,y+b)$.
\begin{enumerate}
  \item Montrer que la translation conserve les distances (si $M\to M'$
        et $N\to N'$, alors $MN=M'N'$).
  \item Montrer que la translation conserve les angles.
  \item Montrer que la translation conserve le parall\'elisme des droites.
  \item Est-ce que la translation conserve les \'equations de droites~?
        Quelle est l'image de $y=x$ par la translation $(1,2)$~?
\end{enumerate}
\end{probleme}

%─────────────────────────────────────────────────────────────────
\subsection*{D\'efis}
%─────────────────────────────────────────────────────────────────

\begin{challenge}[D\'emonstrations g\'eom\'etriques fondamentales]
D\'emontrer par le calcul vectoriel~:
\begin{enumerate}
  \item \textbf{Diagonales d'un losange~:} $ABCD$ est un losange
        ($AB=BC=CD=DA$). Montrer que $\vv{AC}\perp\vv{BD}$.
  \item \textbf{Th\'eor\`eme de la m\'ediane~:} Dans tout triangle $ABC$
        de m\'ediane $AM$ ($M$ milieu de $[BC]$)~:
        $AB^2+AC^2 = 2AM^2+2BM^2$.
  \item \textbf{Angle inscrit dans un demi-cercle~:}
        Si $[AB]$ est un diam\`etre et $M$ un point du cercle,
        alors $\vv{MA}\cdot\vv{MB}=0$.
\end{enumerate}
\end{challenge}

\begin{challenge}[G\'eom\'etrie et coordonn\'ees — droite d'Euler]
\index{droite d'Euler}
Pour tout triangle $ABC$, les points remarquables suivants sont align\'es
sur la \emph{droite d'Euler}~: le circocentre $O'$, le centre de gravité $G$, l'orthocentre $H$.
De plus, $\vv{O'G} = \frac{1}{3}\vv{O'H}$.
\begin{enumerate}
  \item Calculer $O'$, $G$, $H$ pour $A(0,0)$, $B(4,0)$, $C(1,3)$.
  \item V\'erifier l'alignement.
  \item V\'erifier la relation $\vv{O'G} = \frac{1}{3}\vv{O'H}$.
  \item Chercher la signification g\'eom\'etrique du cercle des neuf points
        (cercle passant par les milieux des c\^ot\'es).
\end{enumerate}
\end{challenge}

\begin{challenge}[Produit vectoriel en 2D~: aire orient\'ee]
En dimension $2$, le \emph{d\'eterminant} de deux vecteurs
$\vec{u}=(a,b)$ et $\vec{v}=(c,d)$ est $\det(\vec{u},\vec{v})=ad-bc$.
L'aire du parall\'elogramme form\'e par $\vec{u}$ et $\vec{v}$
vaut $|\det(\vec{u},\vec{v})|$.
\begin{enumerate}
  \item Calculer l'aire du triangle $A(1,2)$, $B(4,6)$, $C(-1,3)$
        en utilisant~: aire $= \frac{1}{2}|\det(\vv{AB},\vv{AC})|$.
  \item D\'emontrer que trois points $A$, $B$, $C$ sont align\'es ssi
        $\det(\vv{AB},\vv{AC})=0$.
  \item Calculer l'aire du quadrilat\`ere $A(0,0)$, $B(3,1)$, $C(4,4)$, $D(1,3)$.
  \item Montrer que $\det(\vec{u}+\vec{v},\vec{w}) = \det(\vec{u},\vec{w})+\det(\vec{v},\vec{w})$.
\end{enumerate}
\end{challenge}

\begin{challenge}[Transformations et vecteurs]
\begin{enumerate}
  \item La sym\'etrie centrale de centre $O$ transforme $\vv{OM}$ en $-\vv{OM}$.
        Trouver l'image de $A(2,3)$, $B(-1,4)$, et du segment $[AB]$
        par la sym\'etrie de centre $O(1,1)$.
  \item La r\'eflexion d'axe $(Ox)$ transforme $(x,y)$ en $(x,-y)$.
        V\'erifier qu'elle conserve les distances.
  \item Composant une translation de vecteur $(1,0)$ et une sym\'etrie centrale
        de centre $O$~: d\'ecrire la transformation r\'esultante.
  \item Montrer qu'une translation suivie d'une autre translation est encore
        une translation. Quelle est son vecteur~?
\end{enumerate}
\end{challenge}

\begin{challenge}[Vecteurs et d\'emonstration analytique]
\begin{enumerate}
  \item Les vecteurs $\vec{u}=(1,2)$ et $\vec{v}=(-2,1)$ sont orthogonaux.
        Montrer que tout vecteur $\vec{w}=(a,b)$ s'\'ecrit
        $\vec{w}=\lambda\vec{u}+\mu\vec{v}$~: trouver $\lambda$ et $\mu$.
  \item G\'en\'eraliser~: si $\vec{u}$ et $\vec{v}$ sont orthogonaux et non nuls,
        tout vecteur du plan s'\'ecrit $\lambda\vec{u}+\mu\vec{v}$.
        Trouver les formules de $\lambda$ et $\mu$ en fonction de
        $\vec{u}\cdot\vec{w}$ et $\vec{v}\cdot\vec{w}$.
  \item Application~: d\'ecomposer $\vec{w}=(5,3)$ sur la base $\vec{u}=(1,0)$
        et $\vec{v}=(0,1)$. Puis sur la base $\vec{u}=(1,1)/\sqrt{2}$
        et $\vec{v}=(1,-1)/\sqrt{2}$.
\end{enumerate}
\end{challenge}

\begin{challenge}[Vecteurs en physique~: trajectoire]
Un projectile est lanc\'e de $O(0,0)$ avec une vitesse initiale
$\vec{v_0} = (v_x, v_y)$ (m/s). Sa position au temps $t$ est~:
$$\vec{r}(t) = (v_x t,\; v_y t - \tfrac{1}{2}gt^2)$$
avec $g = 10$ m/s$^2$.
\begin{enumerate}
  \item Pour $v_x = 30$ et $v_y = 40$ m/s~:
        calculer la position aux instants $t=0, 1, 2, 3, 4$ s.
  \item \`A quel instant le projectile atteint-il sa hauteur maximale~?
  \item Quand retombe-t-il au sol ($y=0$)~? Quelle est la port\'ee~?
  \item \'Ecrire un algorithme Python \computer\ qui trace la trajectoire.
\end{enumerate}
\end{challenge}

%─────────────────────────────────────────────────────────────────
\begin{bilan}
\begin{itemize}
  \item \textbf{Vecteur~:} direction, sens, norme.
        $\vv{AB}=(x_B-x_A, y_B-y_A)$~;
        $\norme{\vv{AB}}=\sqrt{(x_B-x_A)^2+(y_B-y_A)^2}$.
  \item \textbf{Op\'erations~:} $\vec{u}+\vec{v}=(x_1+x_2,y_1+y_2)$~;
        $\lambda\vec{u}=(\lambda x_1,\lambda y_1)$.
  \item \textbf{Produit scalaire~:}
        $\vec{u}\cdot\vec{v}=x_1x_2+y_1y_2=\norme{\vec{u}}\norme{\vec{v}}\cos\theta$.
  \item \textbf{Orthogonalit\'e~:} $\vec{u}\perp\vec{v}\Leftrightarrow\vec{u}\cdot\vec{v}=0$.
  \item \textbf{Colin\'earit\'e~:} $\vec{u}\parallel\vec{v}
        \Leftrightarrow\det(\vec{u},\vec{v})=x_1y_2-x_2y_1=0$.
  \item \textbf{Milieu~:} $I\!\left(\frac{x_A+x_B}{2},\frac{y_A+y_B}{2}\right)$.
        centre de gravité~: $G\!\left(\frac{x_A+x_B+x_C}{3},\frac{y_A+y_B+y_C}{3}\right)$.
  \item \textbf{Chasles~:} $\vv{AB}+\vv{BC}=\vv{AC}$.
\end{itemize}

%─────────────────────────────────────────────────────────────────
\subsection*{Logique \& Algorithmes}
%─────────────────────────────────────────────────────────────────

\begin{exo}[\logiq\ Propositions sur les vecteurs \corrige]{1}
Pour chacune, dire si vraie ou fausse et \'enoncer la n\'egation~:
\begin{enumerate}
  \item $\forall\vec{u},\vec{v}$~: $\vec{u}+\vec{v}=\vec{v}+\vec{u}$.
  \item $\forall\vec{u}$~: $\vec{u}+\vec{0}=\vec{u}$.
  \item $\exists\vec{u}\ne\vec{0}$~: $\vec{u}\cdot\vec{u}=0$.
  \item $\forall\vec{u},\vec{v}$~: $\vec{u}\cdot\vec{v}=0\Rightarrow\vec{u}=\vec{0}$ ou $\vec{v}=\vec{0}$.
\end{enumerate}
\end{exo}

\begin{exo}[\logiq\ Colin\'earit\'e~: condition n\'ecessaire et suffisante \corrige]{1}
\begin{enumerate}
  \item La condition $\det(\vec{u},\vec{v})=0$ est-elle~: CN, CS, ou CNS pour $\vec{u}\parallel\vec{v}$~?
  \item \'Enoncer l'\'equivalence compl\`ete.
  \item Montrer que $\vec{u}\parallel\vec{v}$ et $\vec{v}\parallel\vec{w}\Rightarrow\vec{u}\parallel\vec{w}$
        (transitivit\'e). C'est quel type de raisonnement~?
\end{enumerate}
\end{exo}

\begin{exo}[\logiq\ Raisonnement par l'absurde en g\'eom\'etrie \corrige]{2}
D\'emontrer par l'absurde~:
\begin{enumerate}
  \item Si $\vv{AB}$ et $\vv{AC}$ sont colin\'eaires et $\norme{\vv{AB}}\ne0$,
        alors $A$, $B$, $C$ sont align\'es.
  \item Si les diagonales d'un quadrilat\`ere se coupent en leur milieu,
        alors c'est un parall\'elogramme.
\end{enumerate}
\end{exo}

\begin{exo}[\logiq\ Implication en g\'eom\'etrie \corrige]{2}
Pour chaque proposition, indiquer si c'est une implication, une \'equivalence,
et si la r\'eciproque est vraie~:
\begin{enumerate}
  \item \og $ABCD$ est un carr\'e $\Rightarrow$ $ABCD$ est un losange.\fg
  \item \og $\vec{u}\cdot\vec{v}=0 \Rightarrow \vec{u}\perp\vec{v}$.\fg
  \item \og $\vv{AB}=\vv{DC} \Rightarrow ABCD$ est un parall\'elogramme.\fg
  \item \og $M$ est milieu de $[AB] \Rightarrow \vv{MA}+\vv{MB}=\vec{0}$.\fg
\end{enumerate}
\end{exo}

\begin{exo}[\logiq\ Conditions pour un rectangle \corrige]{2}
Un parall\'elogramme $ABCD$ est un rectangle ssi ses diagonales sont \'egales.
\begin{enumerate}
  \item Formuler cela comme une \'equivalence.
  \item D\'emontrer le sens direct par le calcul vectoriel.
  \item Formuler la contrapos\'ee.
  \item Donner un exemple de parall\'elogramme non rectangle et v\'erifier
        que ses diagonales sont in\'egales.
\end{enumerate}
\end{exo}

\begin{exo}[\logiq\ Quantificateurs et g\'eom\'etrie \corrige]{3}
On suppose $A$ et $B$ distincts.
\begin{enumerate}
  \item Traduire avec des quantificateurs l'affirmation suivante, puis expliquer
        pourquoi elle est fausse : « Il existe un unique point du plan à égale
        distance de $A$ et de $B$. »
  \item Formuler correctement les deux résultats suivants :
        l'ensemble des points équidistants de $A$ et $B$ est la médiatrice de
        $[AB]$; le milieu est l'unique point de $[AB]$ équidistant de $A$ et $B$.
  \item Écrire avec des quantificateurs : « Tout triangle non aplati possède un
        unique centre de gravité. »
  \item Prouver l'existence et l'unicité de ce point en donnant ses coordonnées
        à partir de celles des trois sommets.
\end{enumerate}
\end{exo}

\begin{exo}[\algo\ Algorithme~: tester si $ABCD$ est un parall\'elogramme \corrige]{1}

\begin{algorithm}[H]
\caption{Test de parall\'elogramme}
\KwIn{coordonn\'ees de $A(x_A,y_A)$, $B(x_B,y_B)$, $C(x_C,y_C)$, $D(x_D,y_D)$}
$\vec{u}\leftarrow\vv{AB}=(x_B-x_A,y_B-y_A)$\;
$\vec{v}\leftarrow\vv{DC}=(x_C-x_D,y_C-y_D)$\;
\Si{$\vec{u}=\vec{v}$}{\Afficher \og Parall\'elogramme\fg\;}
\Sinon{\Afficher \og Pas un parall\'elogramme\fg\;}
\end{algorithm}

Tester pour $A(0,0)$, $B(3,1)$, $C(4,3)$, $D(1,2)$.
Puis pour $A(0,0)$, $B(2,0)$, $C(3,2)$, $D(1,3)$.
\end{exo}

\begin{exo}[\algo\ Algorithme~: nature d'un quadrilat\`ere \computer \corrige]{2}
\'Ecrire un algorithme Python qui, \'etant donn\'e quatre points,
d\'etermine si $ABCD$ est~: parall\'elogramme, rectangle, losange, carr\'e, ou aucun.
Utiliser~: cot\'es, diagonales, produits scalaires.
Tester sur plusieurs exemples.
\end{exo}

\begin{exo}[\algo\ Algorithme~: centre de gravité et milieux \corrige]{1}
\'Ecrire un algorithme qui calcule le centre de gravité d'un triangle $ABC$,
les milieux des cot\'es, et v\'erifie que les m\'edianes se croisent en $G$.
Impl\'ementer \computer.
\end{exo}

\begin{exo}[\algo\ Algorithme~: orthocentre d'un triangle \corrige]{2}
\'Ecrire un programme Python \computer\ qui calcule l'orthocentre $H$ d'un triangle
en trouvant l'intersection des hauteurs.
V\'erifier que le troisi\`eme pied de hauteur passe par $H$.
Tester sur $A(0,4)$, $B(-2,0)$, $C(3,0)$.
\end{exo}

\begin{exo}[\algo\ Algorithme~: produit scalaire et angle \computer \corrige]{2}
\'Ecrire un programme qui, \'etant donn\'es trois points $A$, $B$, $C$,
calcule les trois angles du triangle en utilisant le produit scalaire,
v\'erifie que la somme vaut $\pi$, et affiche le type du triangle.
\end{exo}

\begin{exo}[\algo\ Algorithme~: distance d'un point \`a une droite \corrige]{1}
\'Ecrire un algorithme puis l'impl\'ementer \computer\ pour calculer
la distance du point $A(x_0,y_0)$ \`a la droite $ax+by+c=0$.
Tester pour $A(3,4)$ et la droite $3x+4y-25=0$.
\end{exo}

%─────────────────────────────────────────────────────────────────
\subsection*{Probl\`emes Logique \& Algorithmes}
%─────────────────────────────────────────────────────────────────

\begin{probleme}[\logiq\ Logique et preuves g\'eom\'etriques \corrigeP]{2}
D\'emontrer rigoureusement~:
\begin{enumerate}
  \item Les diagonales d'un losange sont perpendiculaires.
        \'Enoncer l'\'equivalence~: \og $ABCD$ est un losange $\Leftrightarrow$
        $AB=BC$ et $\vv{AC}\perp\vv{BD}$\fg\ (admettre le second membre).
  \item La m\'ediatrice d'un segment est l'ensemble des points \'equidistants
        des deux extr\'emit\'es.
        Formuler cela avec des quantificateurs et d\'emontrer.
\end{enumerate}
\end{probleme}

\begin{probleme}[\logiq\ Logique et propri\'et\'es du produit scalaire \corrigeP]{2}
\begin{enumerate}
  \item Montrer que $\vec{u}\cdot\vec{v}=\vec{v}\cdot\vec{u}$
        (commutativit\'e). C'est une identit\'e~: vraie pour tout $\vec{u},\vec{v}$~?
  \item Montrer que $\vec{u}\cdot(\vec{v}+\vec{w})=\vec{u}\cdot\vec{v}+\vec{u}\cdot\vec{w}$
        (distributivit\'e).
  \item La r\'egle \og $\vec{u}\cdot\vec{v}=0 \Rightarrow \vec{u}=\vec{0}$ ou $\vec{v}=\vec{0}$\fg\
        est-elle vraie~? Si non, donner un contre-exemple.
\end{enumerate}
\end{probleme}

\begin{probleme}[\algo\ Algorithme~: cercle circonscrit \corrigeP]{2}
\'Ecrire un programme Python \computer\ qui calcule le cercle circonscrit
\`a un triangle $ABC$~:
\begin{enumerate}
  \item Calculer les m\'ediatrices de deux c\^ot\'es.
  \item Trouver leur intersection (circocentre $O'$).
  \item Calculer le rayon $R=O'A$.
  \item Afficher l'\'equation du cercle.
  \item Tester sur $A(0,0)$, $B(4,0)$, $C(1,3)$.
\end{enumerate}
\end{probleme}

\begin{probleme}[\algo\ Algorithme~: tracer une trajectoire vectorielle \computer \corrigeP]{2}
Un point part de $A(0,0)$ et se d\'eplace selon des vecteurs $\vec{v_1},\ldots,\vec{v_n}$.
\'Ecrire un programme Python \computer\ qui~:
\begin{enumerate}
  \item Lit les vecteurs de d\'eplacement.
  \item Calcule chaque position interm\'ediaire.
  \item Affiche la trajectoire avec \texttt{matplotlib}.
  \item Calcule la distance totale parcourue.
\end{enumerate}
\end{probleme}

%─────────────────────────────────────────────────────────────────
\subsection*{D\'efis Logique \& Algorithmes}
%─────────────────────────────────────────────────────────────────

\begin{challenge}[\logiq\ Droite d'Euler et logique]
La droite d'Euler dit~: pour tout triangle non \'equilat\'eral,
le circocentre $O'$, le centre de gravité $G$ et l'orthocentre $H$ sont align\'es.
\begin{enumerate}
  \item Formuler ce r\'esultat avec des quantificateurs.
  \item Quel est le statut de ce r\'esultat pour un triangle \'equilat\'eral~?
        La proposition \og pour tout triangle\fg\ est-elle vraie ou fausse~?
  \item V\'erifier sur un exemple num\'erique.
  \item \'Enoncer la contrapos\'ee~: si les trois points ne sont pas align\'es,
        que peut-on conclure~?
\end{enumerate}
\end{challenge}

\begin{challenge}[\algo\ Algorithme g\'en\'eral~: polygone r\'egulier \computer]
\'Ecrire un programme Python \computer\ qui~:
\begin{enumerate}
  \item G\'en\`ere les $n$ sommets d'un polygone r\'egulier inscrit
        dans le cercle de rayon $r$ centr\'e en $O$.
  \item Calcule le p\'erim\`etre et l'aire.
  \item V\'erifie que pour $n\to+\infty$, le p\'erim\`etre tend vers $2\pi r$
        en tabulant pour $n=3,4,5,6,8,12,100$.
  \item Trace le polygone avec \texttt{matplotlib}.
\end{enumerate}
\end{challenge}

\end{bilan}


%% ════════════════════════════════════════════════════════════════
